\subsection{Problem and Objective}\label{sec:problem}

\paragraph{Real-life situation.} The university hosts a group of visiting guests (for example, delegates
from a partner school) and plans a one-day city tour for them by service vehicle. The tour leaves the
\textbf{\StartName}, stops once at each of nine landmarks in downtown Cagayan de Oro City, and
returns to the campus at the end of the day. The planners want the order of stops that keeps the total
driving distance as short as possible, which saves fuel and leaves more time at each stop. This is
exactly a TSP.

\paragraph{Area.} Downtown Cagayan de Oro City, Misamis Oriental. All ten locations lie within the
city centre and the neighbouring Poblacion barangays (Lapasan, Licoan, Burgos, Nazareth, Carmen and
Macabalan), in an area of roughly $2.5 \times 2$~km. We restricted the study to downtown so that
every stop is reachable in a few minutes and the whole graph fits on one readable map.

\paragraph{Starting location and required locations.} The start and end point is V0 (\StartName). The nine
required stops are V1--V9, listed with their coordinates in Table~\ref{tab:vertices}
(Section~\ref{sec:method}).

\paragraph{Assumptions.}
\begin{enumerate}
  \item The group travels by car, and between any two stops the driver follows the shortest driving route
        on the public road network.
  \item The distance between two stops is the same in both directions. Where the road data differ
        because of one-way streets, we use the average of the two directions (Section~\ref{sec:edges}).
  \item Traffic, opening hours, waiting and the time spent at each stop are not modelled; they do not
        change the driving distance.
  \item The vehicle has no capacity or fuel limit within one day.
\end{enumerate}

\paragraph{Constraints.} The route must (i) start and end at V0, and (ii) visit every one of V1--V9
\emph{exactly once}. In graph terms, it must be a Hamiltonian cycle of the graph in Section~\ref{sec:method}.

\paragraph{Objective.} Minimise the \textbf{total travel distance in kilometres},
\[
  \min_{\text{tours } T} \; L(T) = \sum_{k=0}^{n-1} w(v_{k}, v_{k+1}), \qquad v_0 = v_n = \mathrm{V0},
\]
where $w(u,v)$ is the road distance between $u$ and $v$ and $n = \NumVertices$ is the number of locations.
